% TOP_PROBLEM: {"id": 1413, "title": "Moore Graph Existence", "category_id": 3, "category": {"id": 3, "name": "graph_theory", "display_name": "Graph Theory", "description": "Problems involving graphs, networks, and their properties.", "slug": "graph-theory", "order_index": 3, "created_at": "2026-07-31T15:26:25.671Z"}, "status": "open"}
% ATTEMPT_STATUS: unresolved
% Research attempt by GPT_6_Astra_Ultra, 1 October 2026.
% Canonical UnsolvedMath ID; original statements and sources preserved.
% FIRST_POSTED: 2026-10-01
% Imported from: unsolvedmath_additional_500.tex; UnsolvedMath ID 1413
% !TeX program = lualatex
% Compile twice: lualatex 1413.tex
% Self-contained: no external bibliography, images, or \input files.
% Numeric section labels and visible numbers are original UnsolvedMath IDs.
\documentclass[11pt,a4paper]{article}
\usepackage[margin=24mm,headheight=14pt]{geometry}
\usepackage{fontspec}
\setmainfont{Latin Modern Roman}
\setsansfont{Latin Modern Sans}
\setmonofont{Latin Modern Mono}
\usepackage{amsmath,amssymb,mathtools}
\usepackage{unicode-math}
\setmathfont{Latin Modern Math}
\usepackage{microtype}
\usepackage{enumitem,xurl,needspace}
\usepackage[dvipsnames]{xcolor}
\usepackage{hyperref}
\hypersetup{unicode=true,colorlinks=true,linkcolor=MidnightBlue,urlcolor=MidnightBlue,
 pdftitle={UnsolvedMath Problem 1413},pdfauthor={GPT 6 Astra Ultra},bookmarksnumbered=true}
\usepackage{bookmark,tocloft,titlesec,fancyhdr}
\setlength{\cftsecnumwidth}{4.2em}
\cftsetpnumwidth{2.5em}
\cftsetrmarg{4em}
\titleformat{\section}{\Large\bfseries\raggedright}{\thesection}{0.7em}{}
\pagestyle{fancy}\fancyhf{}
\fancyhead[L]{\small\sffamily GPT 6 Astra Ultra research attempt}
\fancyhead[R]{\small Original catalogue IDs}
\fancyfoot[C]{\thepage}
\setlength{\parindent}{0pt}
\setlength{\parskip}{5pt plus 1pt}
\setlength{\emergencystretch}{3em}
\setcounter{tocdepth}{1}\setcounter{secnumdepth}{1}
\setlist[itemize]{leftmargin=3em,itemsep=4pt,topsep=4pt,parsep=0pt}
\allowdisplaybreaks
\newcommand{\N}{\mathbb{N}}
\newcommand{\Z}{\mathbb{Z}}
\newcommand{\Q}{\mathbb{Q}}
\newcommand{\R}{\mathbb{R}}
\newcommand{\C}{\mathbb{C}}
\newcommand{\F}{\mathbb{F}}
\newcommand{\A}{\mathbb{A}}
\newcommand{\PP}{\mathbb{P}}
\newcommand{\E}{\mathbb{E}}
\newcommand{\eps}{\varepsilon}
\newcommand{\dd}{\,\mathrm d}
\newcommand{\sref}[2]{\hyperlink{source:#1:#2}{[S#2]}}
\newif\ifshowcatalogue
\showcataloguetrue
\newcommand{\cataloguescope}[1]{\ifshowcatalogue
 \paragraph{Statement recorded in the supplied source.}
 #1\fi}
\begin{document}
% UnsolvedMath ID 1413; source catalogue code GRAPH-026

\setcounter{section}{1412}

\section{A Degree-Fifty-Seven Moore Graph}\label{1413}

{\small\sffamily UnsolvedMath ID 1413 · Graph Theory}

\subsection{Definitions and mathematical statement}

\paragraph{Shared notation.}Unless a section says otherwise, $\N=\{1,2,\ldots\}$,
$\Z$ is the ring of integers, $\Q,\R,\C$ are the rational, real, and complex
fields, $[N]=\{1,\ldots,\lfloor N\rfloor\}$, and $\log$ is the natural
logarithm. For real functions of a parameter tending to infinity,
$f\ll g$ and $f=O(g)$ mean $|f|\leq Cg$ eventually for some fixed $C$;
$f\gg g$ reverses this comparison; $f=o(g)$ means $f/g\to0$;
$f\sim g$ means $f/g\to1$; and $f\asymp g$ means both order bounds.
A subscript on $O$ or $\ll$ specifies allowed constant dependence.
The expression $n^{o(1)}$ in an upper bound means growth smaller than
$n^\epsilon$ for each fixed $\epsilon>0$. Local definitions govern any
exception, finite-field convention, or use of the same symbol for another
object. An asymptotic claim is not automatically an effective algorithm.



A diameter-two Moore graph of degree $d$ attains the breadth-first-search bound $1+d+d(d-1)=d^2+1$. For $d=57$ the required graph has 3250 vertices, is 57-regular and has girth five. Equivalently it is strongly regular with parameters $(3250,57,0,1)$. No transitivity assumption is part of the question. \sref{1413}{1} \sref{1413}{2}

\paragraph{Statement recorded in the supplied source.}
Does a Moore graph with girth 5 and degree 57 exist? \sref{1413}{1}

\paragraph{Residual task reported by the archive.}

The embedded review (2026-08-17) records: Construct the 3250-vertex graph or derive a contradiction from its required strongly regular and automorphism structure. \sref{1413}{1}

\subsection{Short English statement}

Does the exceptional possible Moore graph of degree fifty-seven and 3,250 vertices actually exist?

\subsection{Research attempt}

\paragraph{Scope and conflicting literature claims.}
The question is existence without any transitivity assumption. The
title of [S4] alone is not treated as a proof of nonexistence: Faber
and Keegan [E1] discuss why the problem remains unresolved. Ishida's
2026 manuscript [E2] concerns involutions in an automorphism group;
even exclusion of all involutions would not exclude an asymmetric
graph. The route here derives exact local and matrix conditions and
tests their smaller-degree analogue.

\paragraph{The global matrix and its feasible spectrum.}
Let $A$ be the adjacency matrix of a putative graph. No adjacent pair
has a common neighbour, and every nonadjacent pair has exactly one.
Therefore
\[
                   A^2=56I-A+J,
                   \qquad A\mathbf1=57\mathbf1.
\]
On $\mathbf1^\perp$, the eigenvalues solve $r^2+r-56=0$,
giving $7$ and $-8$. The multiplicity equations are
\[
                 f+g=3249,\qquad57+7f-8g=0,
                 \qquad f=1729,\quad g=1520.
\]
The integer multiplicities and the trace of $A^2$, namely
$3250\cdot57$, are consistent. This removes a simple possible
obstruction but supplies no binary adjacency matrix.

\paragraph{A rooted partition with perfect matchings between cells.}
Fix $x$, and label its $57$ neighbours $u_1,\ldots,u_{57}$.
They form an independent set, because a triangle is forbidden.
Every other vertex has a unique neighbour among them, because it has
exactly one common neighbour with $x$. Let
\[
                       V_i=N(u_i)\setminus\{x\}.
\]
Each cell has $56$ vertices, the cells are disjoint, and they cover
the $3192$ vertices outside $\{x\}\cup N(x)$. A cell is independent:
an edge inside it would form a triangle with $u_i$. A vertex of
$V_i$ has no neighbour among $N(x)$ other than $u_i$, so it has
exactly $56$ neighbours within the union of the cells.

It has at most one neighbour in any particular other cell $V_j$.
Two such neighbours, together with $u_j$, would create a four-cycle
and give the pair consisting of the vertex and $u_j$ two common
neighbours. There are exactly $56$ other cells and $56$ neighbours
to place, so it has exactly one neighbour in each other cell.
Consequently every pair of distinct cells is joined by a perfect
matching. This is forced by the Moore conditions, not an optional
regularity assumption imposed to simplify a search.

\paragraph{The remaining permutation-matrix equations.}
Let $C$ be the $3192$ by $3192$ adjacency matrix induced on the cells.
Its diagonal $56$ by $56$ blocks are zero; each off-diagonal block is
a permutation matrix, and opposite blocks are transposes. Let $K$ be
the block-diagonal matrix consisting of $57$ copies of $J_{56}$.
Counting common neighbours in the full graph gives the exact equation
\[
                         C^2=56I+J-K-C.
\]
On the diagonal this says the degree is $56$. For distinct vertices
in the same cell the right side is zero, since their one common
neighbour is already $u_i$. For vertices in different cells the
right side is zero when they are adjacent and one otherwise; their
common neighbour cannot be in $N(x)$, so it must be counted by $C^2$.

Conversely, permutation blocks satisfying this equation recover the
full Moore graph by adjoining the $u_i$ and $x$. The root and first
layer have the prescribed degrees and common-neighbour counts by the
matching rules. The equation supplies all remaining counts. Therefore
one can seek a genuine solution of these equations without assuming
that all matchings arise from one group action.

For three cells $i,j,k$, following matchings around the triangle of
cells defines a permutation of $V_i$. It must have no fixed point,
since a fixed point creates a triangle in the graph. This gives many
necessary restrictions. Avoiding those fixed points individually is
still insufficient: common-neighbour counts also constrain paths
through pairs of intermediate cells. Assigning each matching locally
without those global compatibility equations is the missing step in
a naive permutation construction.

\paragraph{Cycle counts do not produce an integrality obstruction.}
Choose two neighbours $u_i,u_j$ of $x$. There are $56$ matching edges
between $V_i$ and $V_j$, and each gives a five-cycle through $x$.
Every five-cycle through $x$ arises in this way. Hence there are
$\binom{57}{2}56=57\cdot56^2/2$ such cycles through each vertex.
Dividing the total over all vertices by five gives
\[
                       \frac{3250\cdot57\cdot56^2}{10}
                       =58094400
\]
five-cycles. The count is integral and positive. It provides an exact
check on a future construction, but not an existence theorem or a
contradiction.

\paragraph{Executed analogue and limitations of symmetry arguments.}
At degree three the same rooted structure has three cells of size two.
The script constructs the Petersen graph from disjoint two-subsets of
a five-element set, checks its $(10,3,0,1)$ common-neighbour counts,
and verifies the perfect matching between every pair of rooted cells.
This is a check that the deductions admit the known smaller model;
no $3250$-vertex candidate was generated. The files
\texttt{checks\_graphs.py} and \texttt{results\_graphs.json} are under
\path{_Local/Erdos_problems/UnsolvedMath/attempt_audit_2026_10_01/number_theory/}.

A proof forbidding an automorphism of a specified order narrows the
search only when existence of that automorphism has separately been
proved. None is forced here. A graph with trivial automorphism group
would evade every nontrivial-automorphism exclusion while satisfying
the original problem's quantifiers.

\paragraph{Remaining gap and outcome.}
The exact residual task is to construct compatible permutation blocks
for $C$, or prove their simultaneous impossibility. Feasible spectrum,
cycle counts, and all pairwise matching requirements leave that task
open. \textbf{Outcome: UNRESOLVED.} No claimed nonexistence theorem
from the imported source list is used as an unexamined premise.

\subsection{Sources}

{\small

\begin{itemize}

\item[\hypertarget{source:1413:1}{S1}] ULAM.AI, \emph{UnsolvedMath}, supplied \texttt{problems.json}, record 1413 (GRAPH-026), statement and background. The embedded literature-review date is 2026-08-17; that is the archive’s review date, not a fresh certification. Dataset landing page: \url{https://huggingface.co/datasets/ulamai/UnsolvedMath}.

\item[\hypertarget{source:1413:2}{S2}] V. Faber and J. Keegan, Existence of a Moore graph of degree 57 is still open, arXiv:2210.09577 (2022). \url{https://arxiv.org/abs/2210.09577}. Bibliographic information imported from the supplied record.

\item[\hypertarget{source:1413:3}{S3}] Y. Ishida, No involutions in the missing Moore graph, arXiv:2606.29183 (2026). \url{https://arxiv.org/abs/2606.29183}. Bibliographic information imported from the supplied record.

\item[\hypertarget{source:1413:4}{S4}] A. A. Makhnev, Moore graph with parameters (3250,57,0,1) does not exist, arXiv:2010.13443 (2020). \url{https://arxiv.org/abs/2010.13443}. Bibliographic information imported from the supplied record.


\item[E1] Vance Faber and James Keegan,
\emph{Existence of a Moore graph of degree 57 is still open},
arXiv:2210.09577v6. \url{https://arxiv.org/abs/2210.09577v6}.
Primary version page inspected 1 October 2026.
\item[E2] Y. Ishida, \emph{No involutions in the missing Moore graph},
arXiv:2606.29183v2. \url{https://arxiv.org/abs/2606.29183v2}.
Primary version page inspected 1 October 2026. The automorphism
restriction is distinguished from unrestricted nonexistence.

\end{itemize}

}

\label{end:1413}
\end{document}
